\section{常用公式、定理与结论}

物理公式是解题的工具，以下按模块整理核心公式与常用结论。

\subsection{力学}

\subsubsection{运动学}

\textbf{匀变速直线运动}：
\[
v = v_0 + at,\quad
x = v_0t + \frac12 at^2,\quad
v^2 - v_0^2 = 2ax,\quad
\bar{v} = \frac{v_0+v}{2} = \frac{x}{t}
\]

\textbf{自由落体}（$v_0=0,\;a=g$）：
\[
v = gt,\quad h = \frac12 gt^2,\quad v^2 = 2gh
\]

\textbf{竖直上抛}：上升时间 $t = \dfrac{v_0}{g}$，最大高度 $H = \dfrac{v_0^2}{2g}$，总时间 $T = \dfrac{2v_0}{g}$

\textbf{匀变速推论}：
\begin{itemize}
  \item 相邻相等时间位移差 $\Delta x = aT^2$
  \item 中间时刻速度 $v_{t/2} = \bar{v} = \dfrac{v_0+v}{2}$
  \item 中间位置速度 $v_{x/2} = \sqrt{\dfrac{v_0^2+v^2}{2}}$
  \item 初速为零：$x_1:x_2:x_3:\cdots = 1:3:5:\cdots$
\end{itemize}

\textbf{平抛运动}：
\[
x = v_0t,\quad y = \frac12 gt^2,\quad
v_x = v_0,\quad v_y = gt,\quad
\tan\theta = \frac{v_y}{v_x} = \frac{gt}{v_0}
\]
\[
\text{合位移} = \sqrt{x^2+y^2},\quad
\text{合速度} = \sqrt{v_0^2+(gt)^2}
\]

\textbf{圆周运动}：
\[
v = \omega r,\quad
T = \frac{2\pi}{\omega},\quad
a_n = \frac{v^2}{r} = \omega^2 r,\quad
F_n = m\frac{v^2}{r} = m\omega^2 r
\]

\subsubsection{受力分析}

\textbf{重力}：$G = mg$（$g\approx 9.8\,\text{m/s}^2$，湖南高考常取 $10\,\text{m/s}^2$）

\textbf{弹力}：胡克定律 $F = kx$（$x$ 为形变量）

\textbf{摩擦力}：
\[
f_{\text{静}} \le \mu N,\quad
f_{\text{滑}} = \mu N
\]

\textbf{力的合成与分解}：
\[
F_{\text{合}} = \sqrt{F_1^2+F_2^2+2F_1F_2\cos\theta},\quad
\tan\varphi = \frac{F_2\sin\theta}{F_1+F_2\cos\theta}
\]

\textbf{共点力平衡}：$\sum F_x = 0,\;\sum F_y = 0$

\subsubsection{牛顿运动定律}

\[
F_{\text{合}} = ma
\]

\textbf{连接体问题}：整体法求加速度，隔离法求内力

\textbf{超重与失重}：
\begin{itemize}
  \item 超重：$a$ 向上，$N = m(g+a)$
  \item 失重：$a$ 向下，$N = m(g-a)$
  \item 完全失重：$a=g$，$N=0$
\end{itemize}

\subsubsection{曲线运动与万有引力}

\textbf{运动的合成与分解}：合运动与分运动具有等时性、独立性

\textbf{小船渡河}：
\begin{itemize}
  \item 最短时间：船头垂直河岸，$t_{\min}=d/v_{\text{船}}$
  \item 最短位移：船头指向上游，$\cos\theta = v_{\text{水}}/v_{\text{船}}$
\end{itemize}

\textbf{万有引力定律}：
\[
F = G\frac{Mm}{r^2},\quad G = 6.67\times10^{-11}\,\text{N}\cdot\text{m}^2/\text{kg}^2
\]

\textbf{天体运动}：
\[
G\frac{Mm}{r^2} = m\frac{v^2}{r} = m\omega^2 r = m\frac{4\pi^2}{T^2}r
\]
\[
v = \sqrt{\frac{GM}{r}},\quad
\omega = \sqrt{\frac{GM}{r^3}},\quad
T = 2\pi\sqrt{\frac{r^3}{GM}}
\]

\textbf{黄金代换}：$GM = gR^2$（$R$ 为天体半径，$g$ 为表面重力加速度）

\textbf{三个宇宙速度}：
\begin{itemize}
  \item 第一宇宙速度（环绕速度）：$v_1 = \sqrt{gR} \approx 7.9\,\text{km/s}$
  \item 第二宇宙速度（逃逸速度）：$v_2 \approx 11.2\,\text{km/s}$
  \item 第三宇宙速度：$v_3 \approx 16.7\,\text{km/s}$
\end{itemize}

\textbf{同步卫星}：$T = 24\,\text{h}$，轨道在赤道正上方，$h \approx 3.6\times10^4\,\text{km}$

\subsubsection{功与能}

\textbf{功}：$W = Fx\cos\theta$（$\theta$ 为 $F$ 与 $x$ 夹角）

\textbf{功率}：
\[
P = \frac{W}{t} = Fv\cos\theta,\quad
P_{\text{额}} = Fv_{\text{max}} \quad\text{（机车启动）}
\]

\textbf{动能定理}：
\[
W_{\text{合}} = \Delta E_k = \frac12 mv_2^2 - \frac12 mv_1^2
\]

\textbf{机械能守恒定律}（只有重力或弹力做功）：
\[
E_{k1} + E_{p1} = E_{k2} + E_{p2}
\]
\[
\frac12 mv_1^2 + mgh_1 = \frac12 mv_2^2 + mgh_2
\]

\textbf{功能关系}：
\begin{itemize}
  \item 重力做功 = 重力势能减少量：$W_G = -\Delta E_p$
  \item 弹力做功 = 弹性势能减少量
  \item 合外力做功 = 动能变化量（动能定理）
  \item 除重力（弹力）外其他力做功 = 机械能变化量
\end{itemize}

\subsubsection{动量}

\textbf{动量}：$p = mv$

\textbf{冲量}：$I = Ft$

\textbf{动量定理}：$Ft = \Delta p = mv_2 - mv_1$

\textbf{动量守恒定律}（系统合外力为零）：
\[
m_1v_1 + m_2v_2 = m_1v_1' + m_2v_2'
\]

\textbf{弹性碰撞}（动能守恒）：
\[
v_1' = \frac{(m_1-m_2)v_1 + 2m_2v_2}{m_1+m_2},\quad
v_2' = \frac{(m_2-m_1)v_2 + 2m_1v_1}{m_1+m_2}
\]

\textbf{完全非弹性碰撞}（碰后共速，动能损失最大）：
\[
v_{\text{共}} = \frac{m_1v_1+m_2v_2}{m_1+m_2}
\]

\subsubsection{机械振动与机械波}

\textbf{简谐运动}：
\[
F = -kx,\quad
T = 2\pi\sqrt{\frac{m}{k}}\quad\text{（弹簧振子）}
\]
\[
T = 2\pi\sqrt{\frac{l}{g}}\quad\text{（单摆）}
\]

\textbf{波速公式}：$v = \lambda f = \dfrac{\lambda}{T}$

\textbf{干涉条件}：频率相同、相位差恒定
\begin{itemize}
  \item 加强点：$\Delta d = n\lambda$（$n=0,1,2,\cdots$）
  \item 减弱点：$\Delta d = (2n+1)\dfrac{\lambda}{2}$
\end{itemize}

\textbf{多普勒效应}：波源与观察者靠近时频率升高，远离时频率降低

\subsection{热学（选修 3--3）}

\textbf{分子动理论}：
\[
N_A = 6.02\times10^{23}\,\text{mol}^{-1},\quad
\text{分子直径数量级 }10^{-10}\,\text{m}
\]

\textbf{阿伏加德罗常数应用}：
\[
\text{分子质量 } m_0 = \frac{M}{N_A},\quad
\text{分子数 } n = \frac{m}{M}N_A = \frac{V}{V_m}N_A
\]

\textbf{温度}：$T = t + 273.15$（热力学温标）

\textbf{理想气体状态方程}：
\[
\frac{pV}{T} = C\quad\text{或}\quad
\frac{p_1V_1}{T_1} = \frac{p_2V_2}{T_2}
\]

\textbf{气体实验定律}：
\begin{itemize}
  \item 玻意耳定律（等温）：$p_1V_1 = p_2V_2$
  \item 查理定律（等容）：$\dfrac{p_1}{T_1} = \dfrac{p_2}{T_2}$
  \item 盖-吕萨克定律（等压）：$\dfrac{V_1}{T_1} = \dfrac{V_2}{T_2}$
\end{itemize}

\textbf{热力学第一定律}：
\[
\Delta U = W + Q
\]
\begin{itemize}
  \item 吸热 $Q>0$，放热 $Q<0$
  \item 外界对气体做功 $W>0$，气体对外做功 $W<0$
  \item 内能增加 $\Delta U>0$，内能减少 $\Delta U<0$
\end{itemize}

\textbf{热力学第二定律}：
\begin{itemize}
  \item 不可能从单一热源吸热使之完全变成功而不产生其他影响
  \item 热量不可能自发地从低温物体传到高温物体
\end{itemize}

\textbf{气体压强微观解释}：分子平均动能（温度）$+$ 分子数密度（体积）

\textbf{晶体与非晶体}：
\begin{itemize}
  \item 晶体：有固定熔点，各向异性（单晶）或各向同性（多晶）
  \item 非晶体：无固定熔点，各向同性
\end{itemize}

\textbf{液体表面张力}：使液面收缩，方向沿液面切线方向

\subsection{电学}

\subsubsection{静电场}

\textbf{库仑定律}：
\[
F = k\frac{q_1q_2}{r^2},\quad
k = 9.0\times10^9\,\text{N}\cdot\text{m}^2/\text{C}^2
\]

\textbf{电场强度}：
\[
E = \frac{F}{q},\quad
E = k\frac{Q}{r^2}\;\text{（点电荷）},\quad
E = \frac{U}{d}\;\text{（匀强电场）}
\]

\textbf{电场力做功与电势能}：
\[
W_{AB} = qU_{AB} = -\Delta E_p,\quad
U_{AB} = \frac{W_{AB}}{q} = \varphi_A - \varphi_B
\]

\textbf{电势与电势差}：
\[
\varphi = \frac{E_p}{q},\quad
U_{AB} = Ed\;\text{（匀强电场，$d$ 沿电场线方向）}
\]

\textbf{电容}：
\[
C = \frac{Q}{U},\quad
C = \frac{\varepsilon_r S}{4\pi kd}\;\text{（平行板电容器）}
\]

\textbf{带电粒子在电场中偏转}（类平抛）：
\[
a = \frac{qE}{m} = \frac{qU}{md},\quad
t = \frac{l}{v_0},\quad
y = \frac12 at^2 = \frac{qUl^2}{2mdv_0^2}
\]

\subsubsection{恒定电流}

\textbf{欧姆定律}：
\[
I = \frac{U}{R},\quad
R = \rho\frac{l}{S}\;\text{（电阻定律）}
\]

\textbf{串并联规律}：
\[
R_{\text{串}} = R_1 + R_2 + \cdots,\quad
\frac{1}{R_{\text{并}}} = \frac{1}{R_1} + \frac{1}{R_2} + \cdots
\]

\textbf{电功与电功率}：
\[
W = UIt,\quad
P = UI,\quad
Q = I^2Rt\;\text{（焦耳定律）}
\]

\textbf{闭合电路欧姆定律}：
\[
I = \frac{E}{R+r},\quad
U_{\text{外}} = IR = E - Ir
\]

\textbf{电源功率}：
\[
P_{\text{总}} = EI,\quad
P_{\text{出}} = UI,\quad
P_{\text{内}} = I^2r,\quad
P_{\text{出 max}} = \frac{E^2}{4r}\;(R=r\text{ 时})
\]

\textbf{电表改装}：
\begin{itemize}
  \item 电压表：$R_{\text{串}} = \dfrac{U}{I_g} - R_g$
  \item 电流表：$R_{\text{并}} = \dfrac{I_gR_g}{I - I_g}$
\end{itemize}

\subsubsection{磁场}

\textbf{磁感应强度}：
\[
B = \frac{F}{IL}\quad\text{（通电直导线）}
\]

\textbf{安培力}：
\[
F = BIL\sin\theta\quad\text{（$\theta$ 为 $B$ 与 $I$ 夹角）}
\]

\textbf{洛伦兹力}：
\[
f = qvB\sin\theta,\quad
f = qvB\;(v\perp B)
\]

\textbf{带电粒子在匀强磁场中匀速圆周运动}：
\[
qvB = m\frac{v^2}{r},\quad
r = \frac{mv}{qB},\quad
T = \frac{2\pi m}{qB}
\]

\textbf{回旋加速器}：$E_{k\max} = \dfrac{q^2B^2R^2}{2m}$

\textbf{质谱仪}：$r = \dfrac{1}{B}\sqrt{\dfrac{2mU}{q}}$

\textbf{速度选择器}：$v = \dfrac{E}{B}$（粒子匀速通过）

\subsubsection{电磁感应}

\textbf{法拉第电磁感应定律}：
\[
\mathcal{E} = n\frac{\Delta\Phi}{\Delta t},\quad
\mathcal{E} = Blv\;\text{（导体切割磁感线）}
\]

\textbf{楞次定律}：感应电流的磁场总要阻碍引起感应电流的磁通量的变化

\textbf{自感电动势}：
\[
\mathcal{E}_L = L\frac{\Delta I}{\Delta t}
\]

\textbf{交变电流}：
\[
e = E_m\sin\omega t,\quad
E_m = nBS\omega,\quad
E_{\text{有效}} = \frac{E_m}{\sqrt{2}}
\]
\[
u = U_m\sin\omega t,\quad
i = I_m\sin(\omega t+\varphi)
\]

\textbf{变压器}：
\[
\frac{U_1}{U_2} = \frac{n_1}{n_2},\quad
\frac{I_1}{I_2} = \frac{n_2}{n_1}\;\text{（单相理想变压器）}
\]

\textbf{远距离输电}：$P_{\text{损}} = I^2R = \dfrac{P^2}{U^2}R$

\subsubsection{电磁振荡与电磁波}

\[
T = 2\pi\sqrt{LC},\quad
f = \frac{1}{2\pi\sqrt{LC}}
\]

\textbf{麦克斯韦电磁场理论}：变化的电场产生磁场，变化的磁场产生电场

\textbf{电磁波谱}（频率从低到高）：无线电波 $\to$ 微波 $\to$ 红外线 $\to$ 可见光 $\to$ 紫外线 $\to$ X 射线 $\to$ $\gamma$ 射线

\subsection{光学}

\textbf{折射定律}：
\[
n = \frac{\sin i}{\sin r} = \frac{c}{v}
\]

\textbf{全反射}：
\[
\sin C = \frac{1}{n}\quad\text{（$C$ 为临界角）}
\]

\textbf{光速}：$c = 3.0\times10^8\,\text{m/s}$

\textbf{干涉}：
\begin{itemize}
  \item 双缝干涉：$\Delta x = \dfrac{l}{d}\lambda$（亮纹中心距）
  \item 亮纹条件：$\delta = k\lambda$（$k=0,1,2,\cdots$）
  \item 暗纹条件：$\delta = (2k+1)\dfrac{\lambda}{2}$
\end{itemize}

\textbf{衍射}：光绕过障碍物的现象，发生明显衍射的条件：障碍物尺寸与波长可比

\textbf{偏振}：证明光是横波

\textbf{光电效应}：
\[
E_k = h\nu - W_0,\quad
h = 6.63\times10^{-34}\,\text{J}\cdot\text{s}
\]
\[
\nu_{\text{截止}} = \frac{W_0}{h},\quad
\text{遏止电压 } U_c = \frac{h\nu - W_0}{e}
\]

\textbf{爱因斯坦光子说}：$E = h\nu$

\textbf{光的波粒二象性}：干涉、衍射体现波动性；光电效应体现粒子性

\subsection{原子物理}

\textbf{玻尔原子模型}：
\[
h\nu = E_m - E_n\;(m>n),\quad
r_n = n^2 r_1,\quad
E_n = \frac{E_1}{n^2}\;(E_1=-13.6\,\text{eV})
\]

\textbf{氢原子能级图}：
\[
n=1:\;-13.6\,\text{eV},\quad
n=2:\;-3.4\,\text{eV},\quad
n=3:\;-1.51\,\text{eV},\quad
n=4:\;-0.85\,\text{eV},\;\cdots
\]

\textbf{原子核衰变}：
\begin{itemize}
  \item $\alpha$ 衰变：$\prescript{A}{Z}{X} \to \prescript{A-4}{Z-2}{Y} + \prescript{4}{2}{\text{He}}$
  \item $\beta$ 衰变：$\prescript{A}{Z}{X} \to \prescript{A}{Z+1}{Y} + \prescript{0}{-1}{e}$
  \item $\gamma$ 衰变：伴随 $\alpha$ 或 $\beta$ 衰变放出 $\gamma$ 光子
\end{itemize}

\textbf{半衰期}：
\[
N = N_0\left(\frac12\right)^{t/\tau},\quad
m = m_0\left(\frac12\right)^{t/\tau}
\]

\textbf{核能}：
\[
\Delta E = \Delta m c^2 = 931.5\,\Delta m\,\text{MeV}\quad\text{（$\Delta m$ 以 u 为单位）}
\]

\textbf{核反应类型}：
\begin{itemize}
  \item 衰变：自发进行
  \item 人工转变：用粒子轰击原子核
  \item 裂变：重核分裂（如 $\prescript{235}{92}{U} + \prescript{1}{0}{n} \to \prescript{144}{56}{Ba} + \prescript{89}{36}{Kr} + 3\prescript{1}{0}{n}$）
  \item 聚变：轻核结合（如 $\prescript{2}{1}{H} + \prescript{3}{1}{H} \to \prescript{4}{2}{He} + \prescript{1}{0}{n}$）
\end{itemize}

\textbf{质量亏损}：原子核的质量小于组成它的核子质量之和

\subsection{实验常用公式}

\textbf{纸带处理}：
\[
v_n = \frac{x_n + x_{n+1}}{2T},\quad
a = \frac{x_4+x_5+x_6 - (x_1+x_2+x_3)}{9T^2}
\]

\textbf{伏安法测电阻}：
\[
\text{内接法：} R_{\text{测}} = R_A + R_x > R_x,\quad
\text{外接法：} R_{\text{测}} = \frac{R_xR_V}{R_x+R_V} < R_x
\]

\textbf{选择原则}：$\dfrac{R_x}{R_A} > \dfrac{R_V}{R_x}$ 时用内接法（大电阻），反之用外接法

\textbf{多用电表读数}：欧姆表读数 $=$ 指针示数 $\times$ 倍率

\textbf{等效法测电源电动势和内阻}：
\[
E = \frac{I_1U_2 - I_2U_1}{I_1 - I_2},\quad
r = \frac{U_2 - U_1}{I_1 - I_2}
\]

\subsection{常用物理常量速查}

\[
\begin{array}{c|c}
\text{常量} & \text{数值} \\ \hline
\text{重力加速度 } g & 9.8\;\text{m/s}^2 \\
\text{万有引力常量 } G & 6.67\times10^{-11}\;\text{N}\cdot\text{m}^2/\text{kg}^2 \\
\text{真空光速 } c & 3.0\times10^8\;\text{m/s} \\
\text{元电荷 } e & 1.60\times10^{-19}\;\text{C} \\
\text{电子质量 } m_e & 9.11\times10^{-31}\;\text{kg} \\
\text{质子质量 } m_p & 1.673\times10^{-27}\;\text{kg} \\
\text{中子质量 } m_n & 1.675\times10^{-27}\;\text{kg} \\
\text{普朗克常量 } h & 6.63\times10^{-34}\;\text{J}\cdot\text{s} \\
\text{静电力常量 } k & 9.0\times10^9\;\text{N}\cdot\text{m}^2/\text{C}^2 \\
\text{阿伏加德罗常数 } N_A & 6.02\times10^{23}\;\text{mol}^{-1} \\
\text{标准大气压 } p_0 & 1.013\times10^5\;\text{Pa} \\
\text{热力学零度} & -273.15\;^\circ\text{C}
\end{array}
\]
